\section{Introduction}

EEG research in ADHD has reported heterogeneous task-related oscillatory findings across theta, alpha, beta, and event-related measures \parencite{michelini2022eventoscillations,arpaia2022eegfeatures}. This heterogeneity is scientifically important because a wearable device does not become a valid state monitor merely by reproducing a familiar frequency band.

Wireless and wearable systems offer practical advantages in pediatric and ecologically oriented research. A recent scoping review of wireless EEG in children with developmental disabilities found broad use across ADHD and other populations, but also identified incomplete reporting of data-quality procedures and limited direct validation against reference systems \parencite{park2026wirelesseeg}. Wearable EEG additionally faces artifact-control constraints associated with motion, reduced channel coverage, and dry or semi-wet electrodes \parencite{wearableartifacts2025}.

The present review therefore does not ask whether EEG can diagnose ADHD. Its primary question is narrower: under what measurement and reporting conditions can wearable or low-density EEG support reproducible task-locked or intervention-locked inference in ADHD and related pediatric neurodevelopmental research? The focus is on timing, signal quality, artifacts, task context, repeatability, and evidence boundaries needed for mechanistic monitoring.
